\scanpage{33}
\begin{center}
Cours DEA 1977/78 de A. Grothendieck

\uline{Géométrie arithmétique des polygônes réguliers}
\end{center}

{\reversemarginpar\margin{\handwritten{!}}\normalmarginpar}%
% The hand of the marginal ! and short underline is unassigned.
% The short hand-drawn rule underlines the typed words en commun.
Ce ``cours'' se fera dans l'esprit d'une investigation \handwritten{\uline}{en commun} de multiples facettes de la notion de polygône régulier au-dessus d'un corps de base, puis d'un anneau de base quelconque. Ce sera une occasion stimulante pour nous familiariser avec les notions de base de géométrie vectorielle, \tikz[remember picture,overlay]{\coordinate (s033C-affine);}affine, projective - avec ou sans métrique - en dimensions un, deux et trois, et avec les phénomènes particuliers à ces dimensions, ainsi qu'avec les aspects parfois imprévus que prennent des notions familières (comme celle de réflexion\deleted{s}, ou de centre d'un polygône régulier) sur les corps de caractéristique positive, et notamment sur les corps finis. Le besoin d'une géométrie sur des anneaux de base généraux (tel l'anneau ``sur lequel s'écrit'' le $\mathrm{n}$-polygône régulier ``le plus général'') se fera sentir en cours de route, et nous amènera sans doute à affleurer quelques points de fondements de géométrie algébrique. Si le temps le permet et si d'autres investigations ne nous retiennent , nous tenterons, dans la volée, de construire l'icosaèdre régulier ``général'' sur un anneau de base convenable, fini sur l'anneau ${\handwritten{\mathbb}{Z}}$ des entiers, ayant des corps résiduels de toutes caractéristiques. (La caractéristique deux \hspace{.5em}donnera sans doute du fil à retordre !).

A partir d'un point de départ intuitif et familier, ce ``cours'' voudrait\tikz[remember picture,overlay]{\coordinate (s033C-voudrait);} être une ouverture à un univers fascinant, pleins de mystères qui ne demandent qu'à être élucidés. Il ne peut donner des fruits qu'avec une participation \tikz[remember picture,overlay]{\coordinate (s033C-active);}active de l'étudiant et un travail personnel important, notamment pour développer a\deleted{\unsure{o|e}}vec\\
\tikz[remember picture,overlay]{\coordinate (s033C-detail);}le détail nécessaire les nombreux points \handwritten{(}de fondement\inserted{\handwritten{s}} notamment) qui n'auront été qu'esquissés dans le ``cours''. Il lui est très fortement conseillé de suivre en même temps le cours C4 sur l'icosaèdre, qui se fera dans un esprit analogue. Il y aura sans doute de nombreuses interférences entre ces deux cours de réflexion.

% Detached geometries: no verbal role, origin or author is assigned.
% Local source-pixel geometry is scaled to the typeset x-height; the zero-width
% text anchors retain each mark's relation to the corresponding passage.
% The line break before le detail retains the left-edge mark's line attachment.
\figure{\begin{tikzpicture}[remember picture,overlay]
% Pale oval/U-shaped trace below the a of affine, full-source area x=1255--1264, y=638--650.
\begin{scope}[shift={(s033C-affine)},x=0.26pt,y=-0.26pt]
\draw[black!30,line width=0.65pt,line cap=round] (3,0) .. controls (1,4) and (2,9) .. (6,10) .. controls (9,9) and (10,5) .. (8,2);
\end{scope}
% Detached dark hook below the end of voudrait, around x=1498, y=1582.
\begin{scope}[shift={(s033C-voudrait)},x=0.26pt,y=-0.26pt]
\draw[line width=0.75pt,line cap=round] (-17,16) .. controls (-15,20) and (-19,22) .. (-21,22);
\end{scope}
% Detached upper-right / lower-left pair below active, around (1534,1722), (1529,1729).
\begin{scope}[shift={(s033C-active)},x=0.26pt,y=-0.26pt]
\fill (31,16) ellipse[x radius=1.5,y radius=1.1,rotate=45];
\fill (26,23) ellipse[x radius=1.3,y radius=1.0,rotate=20];
\end{scope}
% Compact two-lobed mark at the physical left edge beside le detail, around x=28--37,y=1839--1846.
\begin{scope}[shift={(s033C-detail)},x=0.26pt,y=-0.26pt]
\fill (-227,-5) .. controls (-227,-8) and (-223,-8) .. (-222,-5)
  .. controls (-219,-6) and (-217,-5) .. (-217,-3)
  .. controls (-217,0) and (-220,1) .. (-221,0)
  .. controls (-223,0) and (-223,-3) .. (-224,-3)
  .. controls (-226,-3) and (-228,-3) .. (-227,-5) -- cycle;
\end{scope}
\end{tikzpicture}}
% archive sheet number: 32
